\section{Moving a signal to a larger scale}\label{sec:multiscale}

The signal in \eqref{eq:signal} discounts positions beyond distance about
$1/s$ from the first disagreement.  The next proposition moves such a signal
to a much smaller decay parameter.  It permits an increase in feature order,
but keeps the loss in the logarithm of the signal below a factor of eight.
Controlling the angles of the new parameters is equally important: those
bounds will later permit passage from input coordinates to trace coordinates.
Enlarging spatial scales while increasing statistic order follows the
strategy of Burudgunte, Valiant, and Wang~\cite{BVW2026}.  Here the comparison
functional may be a positive truncated moment functional, and the scale step
uses Gram matrices of polynomial vectors.

\begin{proposition}[One multiscale step]\label{prop:multiscale}
There are absolute constants $0<\beta,\eta_0<1$ and $C_0,C_1\ge1$ with the
following property.  Let $f$ be an order-$k$ feature from
\eqref{eq:feature}, and let $\Delta$ arise from a moment functional of order
$2D$ with the prefix conditions of Section~\ref{sec:setup}.  Suppose
\[
 0<S\le4,\qquad 0<S'\le\min\{1,\eta_0S\},\qquad 4k\le D,
\]
and, for some $\beta S\le s\le S$ and $\abs{\omega}\le2.002s$,
\begin{equation}\label{eq:multiscale-cost}
 \abs{\mathcal S_f(s,\omega)}\ge e^{-E},\qquad
 E\ge C_0\left(1+k\log(2k)+\log H_f+\log\frac1{S'}\right).
\end{equation}
Then an order-$k'$ feature $f'$, with $1\le k'\le4k$, satisfies
\begin{equation}\label{eq:multiscale-output}
 \abs{\mathcal S_{f'}(s',\omega')}\ge e^{-7.98E},\qquad
 \beta S'\le s'\le S',\qquad \abs{\omega'}\le2.002s',
\end{equation}
and
\begin{equation}\label{eq:multiscale-height}
 H_{f'}\le H_f^4\left(1+\frac1{\beta S'}\right)^3.
\end{equation}
Moreover, there are $1\le m\le4$ lists, each equal to the old parameter list
$(z_1,\ldots,z_{k-1})$ or its complex conjugate, and at most $m-1$ added
parameters, with the following property.  Each new inner-gap parameter is a
nonempty product using at most one entry from each list and at most one added
parameter.  Every added parameter $e^{-a+\ii\theta}$ obeys
\begin{equation}\label{eq:added-parameters}
 a\ge\beta S',\qquad \abs{\theta}\le C_1S.
\end{equation}
\end{proposition}

The proof has two parts.  First, a broad Gaussian window turns the original
signal into Fourier sums.  Positivity of the moment functional then forces a
product of at most four such sums to retain a signal while its total frequency
is small.  Second, expanding that product and deforming Gaussian contours
produces a single feature with the asserted parameters.  The first part
occupies the present subsections; the conversion is carried out in
Section~\ref{subsec:contour-conversion}.

\subsection{A Gaussian window and its logarithmic bound}

Throughout the proof of Proposition~\ref{prop:multiscale}, set
\begin{equation}\label{eq:window-parameters}
 \sigma=\frac{K\sqrt E}{S'},\qquad
 W(l)=\exp\left(-\frac{l^2}{2\sigma^2}\right),\qquad b=\beta S',
\end{equation}
where $K$ is a sufficiently large absolute constant.  The constants will be
chosen in the order $K$, then $\beta$ sufficiently small, then $\eta_0$
sufficiently small, and finally $C_0$ sufficiently large.  The choices of $K$
and $\beta$ will also enforce the contour estimates below.

We use the following consequence of \eqref{eq:multiscale-cost} repeatedly.
The logarithm of any fixed product of powers of
$H_f,1+\sigma,1/S',1/b$, together with a factor
$\exp(O(k\log(2k)))$, is bounded by
\[
 C\left(1+k\log(2k)+\log H_f+\log\frac1{S'}\right)+C\log E,
\]
where $C$ depends only on constants already chosen.  Increasing $C_0$ makes
this at most any prescribed positive multiple of $E$.  Only finitely many
such estimates will be required, so a single final choice of $C_0$ suffices.

For real $t$, define the polynomial and its evaluation
\begin{equation}\label{eq:window-sums}
 A(t)=\sum_{j=1}^n f_jW(j-d)e^{\ii t(j-d)},\qquad B_x(t)=A(t)(x).
\end{equation}
The coefficient bound for $f_j$ and the Gaussian sum imply that
\begin{equation}\label{eq:window-uniform-bound}
 \abs{B_x(t)},\quad \abs{\EE[A(t)]},\quad \abs{\Delta[A(t)]}
 \le U:=C H_f(1+\sigma),
\end{equation}
with $C$ absolute and chosen so that $U\ge1$.
By prefix cancellation,
\[
 P(z):=\sum_{j=d}^n\Delta[f_j]W(j-d)z^{j-d}
 \quad\hbox{satisfies}\quad P(e^{\ii t})=\Delta[A(t)].
\]
We first show that $P$ retains a signal at almost the original decay:
\begin{equation}\label{eq:window-interior-signal}
 \abs{P(e^{-s_1+\ii\omega})}\ge\tfrac12e^{-E}
 \quad\hbox{for some}\quad \abs{s_1/s-1}\le10^{-4}.
\end{equation}

Let $V$ be a centered real normal variable with variance $\sigma^{-2}$.
Since $W(l)\mathbf E e^{Vl}=1$,
\[
 \mathcal S_f(s,\omega)=\mathbf E\,P(e^{-s+V+\ii\omega}).
\]
Write $\delta=10^{-4}$.  Exponential tilting gives
\[
 W(l)\mathbf E\big[e^{Vl}\mathbf1_{\{\abs V>\delta s\}}\big]
 =\Pr\!\left(\abs{N(l/\sigma^2,\sigma^{-2})}>\delta s\right).
\]
For $0\le l\le\delta s\sigma^2/2$, this probability is at most
$2e^{-c\delta^2s^2\sigma^2}$.  For larger $l$, use the bound one and sum
$e^{-sl}$.  Because $\abs{\Delta[f_j]}\le2H_f$, the contribution from
$\abs V>\delta s$ is at most
\begin{equation}\label{eq:window-tail}
 2H_f(1+1/s)\left(2e^{-c\delta^2s^2\sigma^2}
                      +e^{-\delta s^2\sigma^2/2}\right).
\end{equation}
Here $1/s\le1/(\beta S')$, and
$s^2\sigma^2\ge\beta^2K^2E/\eta_0^2$.  Choose $\eta_0$ sufficiently small
and then $C_0$ sufficiently large to make \eqref{eq:window-tail} less than
$e^{-2E}$.  The integral over $\abs V\le\delta s$ consequently has magnitude
at least $e^{-E}-e^{-2E}\ge e^{-E}/2$.  Its maximum integrand magnitude is
at least that large, proving \eqref{eq:window-interior-signal}.

Set
\[
 \Lambda(t)=\log\frac{U}{\abs{P(e^{\ii t})}}\ge0,
\]
with value $+\infty$ at a boundary zero.  We claim
\begin{equation}\label{eq:poisson-log-bound}
 \int_{\RR}\Lambda(t)\frac{s_1\,dt}
 {\pi(s_1^2+(t-\omega)^2)}\le1.001E.
\end{equation}
For completeness, the Cauchy Fourier transform is
\begin{equation}\label{eq:cauchy-transform}
 \int_{\RR}\frac{e^{\ii\xi v}}{\pi(1+v^2)}\,dv=e^{-\abs\xi}.
\end{equation}
For $\xi>0$, close the contour in the upper half-plane and take the residue
at $\ii$; for $\xi<0$ use the lower half-plane, and for $\xi=0$ integrate
directly.  Thus periodizing the density in
\eqref{eq:poisson-log-bound} gives the circle Poisson measure at
$e^{-s_1+\ii\omega}$.

To apply it, factor the nonzero polynomial $P$ into linear factors.  A factor
with a root outside the disk has harmonic logarithmic magnitude inside.  For
a root $a$ inside, the boundary identity
$\abs{e^{\ii t}-a}=\abs{1-\overline a e^{\ii t}}$ and the inequality
\[
 \abs{1-\overline a z}^2-\abs{z-a}^2
 =(1-\abs a^2)(1-\abs z^2)\ge0
\]
give the corresponding upper bound by its Poisson integral.  Roots on the
circle follow by a radial limit; their logarithmic singularities are
integrable.  Summing over factors yields
\[
 \int_{\RR}\Lambda(t)\frac{s_1\,dt}{\pi(s_1^2+(t-\omega)^2)}
 \le\log U-\log\abs{P(e^{-s_1+\ii\omega})}
 \le E+\log(2U).
\]
The last term is at most $1.001E$ by the choice of $C_0$, proving the claim.

\subsection{Frequencies at prescribed differences}

We seek a frequency $\lambda\in[-2s_1,2s_1]$ with a large signal and pairs
of frequencies with every prescribed difference in that interval.
The product $A(t)\overline{A(t')}$ has signed frequency $t-t'$.
After normalization by their evaluations at $x$, these products will supply
the interior vectors of a chain from $1$ to a normalized copy of
$A(\lambda)$, with consecutive frequency labels differing by at most $b$.
Each neighboring inner product will then have
signed frequency of magnitude at most $b$, even though $\abs\lambda$ may
be much larger.  The logarithmic estimate makes the required existence of
pairs a question about pairing Cauchy mass.

\begin{lemma}[Cauchy matching]\label{lem:cauchy-matching}
Let $w(v)=1/(\pi(1+v^2))$.  For every $r\ge0$ there is a finite positive
measure $\mu_r$ on $\RR^2$, supported on $v-v'=r$, such that the sum of its
two marginals is bounded above by $w(v)\,dv$.  For $r>0$ its total mass is
\begin{equation}\label{eq:cauchy-matching-mass}
 2\mu_r(\RR^2)=1-\frac4\pi\arctan(e^{-\pi/r});
\end{equation}
for $r=0$ its total mass is $1/2$.  In particular, the quantity in
\eqref{eq:cauchy-matching-mass} is greater than $0.737$ when $0<r\le2$.
\end{lemma}

\begin{proof}
For $r=0$, put half the measure $w(v)\,dv$ on the diagonal.  Suppose $r>0$,
and partition the line into chains $v+jr$, $j\in\mathbb Z$, with
$-r/2\le v<r/2$.  Abbreviate $w_j=w(v+jr)$.  On the edge from $j$ to $j-1$,
for $j\ge1$, put weight
\[
 e_j=\sum_{m\ge0}(-1)^m w_{j+m}.
\]
These weights are nonnegative, since the summands decrease along the tail.
Moreover $e_j+e_{j+1}=w_j$, so the edges exhaust the weight at every positive
vertex.  Construct the negative-side edges in the same way, proceeding from
the negative tail toward zero.  The only possibly unused weight is at the
central vertex, where its value is
\[
 R(v)=\sum_{j\in\mathbb Z}(-1)^j w(v+jr).
\]
To check that no weight has been overused, periodize $w$ at period $2r$ and
subtract its translate by $r$.  By \eqref{eq:cauchy-transform}, with
$\tau=e^{-\pi/r}$, this gives the absolutely convergent Fourier series
\[
 R(v)=\frac2r\sum_{m\ge0}\tau^{2m+1}
                  \cos\frac{(2m+1)\pi v}{r}.
\]
Writing $a=\pi v/r$, the sum equals
\[
 \re\frac{\tau e^{\ii a}}{1-\tau^2e^{2\ii a}}
 =\frac{\tau(1-\tau^2)\cos a}
        {1-2\tau^2\cos(2a)+\tau^4}\ge0,
\]
since $\abs a\le\pi/2$.  Orient every edge with its larger endpoint first
and integrate the edge weights with respect to $dv$ over $[-r/2,r/2)$.
The resulting measure has the asserted marginal domination.  Its uncovered
mass is
\[
 \int_{-r/2}^{r/2}R(v)\,dv
 =\frac4\pi\sum_{m\ge0}\frac{(-1)^m\tau^{2m+1}}{2m+1}
 =\frac4\pi\arctan\tau,
\]
which proves \eqref{eq:cauchy-matching-mass}.  The expression decreases with
$r$, and at $r=2$ it exceeds $0.737$.  For example,
$e^{-\pi/2}<0.209$ together with
$\arctan u\le u-u^3/3+u^5/5$ for $0\le u\le1$ proves this numerical bound.
\end{proof}

\begin{lemma}[Frequencies from the logarithmic bound]\label{lem:good-frequencies}
Let $s_1,E>0$, let $\abs\omega\le2.003s_1$, and let
$\Lambda:\RR\to[0,\infty]$ be measurable and satisfy
\eqref{eq:poisson-log-bound}.  There is a real $\lambda$ with
\begin{equation}\label{eq:single-frequency}
 \abs\lambda\le2s_1,\qquad \Lambda(\lambda)\le2.40E.
\end{equation}
There is also an absolute $C_2$ such that, for every $\abs y\le2s_1$, some
real $t,t'$ satisfy
\begin{equation}\label{eq:paired-frequencies}
 t-t'=y,\qquad \abs t,\abs{t'}\le C_2s_1,\qquad
 \Lambda(t)+\Lambda(t')\le2.74E.
\end{equation}
\end{lemma}

\begin{proof}
The Cauchy mass of $[-2s_1,2s_1]$ is at least
\[
 \frac{\arctan(2-2.003)+\arctan(2+2.003)}\pi>0.418.
\]
Because $1.001/0.418<2.40$, averaging proves
\eqref{eq:single-frequency}.

For \eqref{eq:paired-frequencies}, use Lemma~\ref{lem:cauchy-matching} with
$r=\abs y/s_1\le2$ and then substitute $t=\omega+s_1v$.  Reverse the
orientation if needed.  The sum of the two marginals is bounded by the
Cauchy measure in \eqref{eq:poisson-log-bound}, while twice the total pair
mass is greater than $0.737$.  Discard pairs with an endpoint outside
$[-C_2s_1,C_2s_1]$.  Twice the discarded pair mass is at most twice the
Cauchy tail outside that interval, which tends uniformly to zero as
$C_2\to\infty$, since $\abs\omega/s_1\le2.003$.
Choose an absolute $C_2$ so the retained pair mass exceeds $0.733/2$.
Its integral of $\Lambda(t)+\Lambda(t')$ is at most $1.001E$.  Some retained
pair consequently has sum at most
\[
 \frac{2(1.001)}{0.733}E<2.74E,
\]
as required.  The diagonal construction gives the same argument when $y=0$.
\end{proof}

\subsection{A product with small total frequency}

The window parameters satisfy $\abs\omega/s_1\le2.002/0.9999<2.003$, so
Lemma~\ref{lem:good-frequencies} applies to our $\Lambda$.
Individual frequencies in \eqref{eq:paired-frequencies} need not be small.
Their differences, however, can be prescribed.  We use these pairs to build
a chain of polynomial vectors whose neighboring inner products have small
total frequency.  The following Hilbert-space fact explains why such a chain
preserves the needed exponential accuracy.

\begin{lemma}[Gram chaining]\label{lem:gram-chain}
Let $v_0,\ldots,v_J$, with $J\ge1$, lie in a complex Hilbert space whose
inner product is linear in its first argument.  Suppose $\norm{v_0}=1$ and,
for some $\xi\ge0$,
\[
 \abs{\norm{v_i}^2-1}\le\xi\quad(0\le i\le J),\qquad
 \abs{\langle v_i,v_{i-1}\rangle-1}\le\xi\quad(1\le i\le J).
\]
Then
\[
 \abs{\langle v_J,v_0\rangle-1}\le2J^2\xi.
\]
\end{lemma}

\begin{proof}
Put $h_i=v_i-v_{i-1}$.  The hypotheses give
\[
 \norm{h_i}\le2\sqrt\xi,\qquad
 \abs{\langle h_i,v_{i-1}\rangle}\le2\xi,
 \qquad \norm{v_{i-1}-v_0}\le2(i-1)\sqrt\xi.
\]
Thus Cauchy--Schwarz implies
\[
 \abs{\langle h_i,v_0\rangle}
 \le2\xi+4(i-1)\xi.
\]
Sum over $i$ and use $\langle v_0,v_0\rangle=1$.
\end{proof}

Choose an absolute $C_3$ large enough that all frequencies supplied by
Lemma~\ref{lem:good-frequencies} have magnitude at most $C_3S$.
A product of one to four factors, each of the form $A(t)$ or
$\overline{A(t)}$ with $\abs t\le C_3S$, is called \emph{balanced} if the
sum of its signed frequencies has magnitude at most $b$; conjugation changes
the sign of a frequency.

\begin{lemma}[A balanced product retains a signal]\label{lem:balanced-product}
Some balanced product $G$ satisfies
\begin{equation}\label{eq:balanced-product-signal}
 \abs{\Delta[G]}\ge e^{-7.94E}.
\end{equation}
\end{lemma}

\begin{proof}
Suppose to the contrary that every balanced product has signal less than
$\epsilon=e^{-7.94E}$.  The squared modulus $\abs{A(t)}^2$ is one such
product.  Positivity of $\EE$, applied to $A(t)$ and $1$, gives
\[
 \abs{\EE[A(t)]}^2\le\EE[\abs{A(t)}^2]
 \le\abs{B_x(t)}^2+\epsilon,
\]
and hence
\begin{equation}\label{eq:evaluation-lower-bound}
 \abs{\Delta[A(t)]}\le2\abs{B_x(t)}+\sqrt\epsilon.
\end{equation}
At $\lambda$ from \eqref{eq:single-frequency}, this implies
\begin{equation}\label{eq:endpoint-normalization}
 \abs{B_x(\lambda)}\ge e^{-2.41E}.
\end{equation}
Indeed $\abs{\Delta[A(\lambda)]}\ge e^{-2.40E}$, whereas
$\sqrt\epsilon=e^{-3.97E}$; fixed numerical factors are absorbed by the
absolute lower bound on $E$ in \eqref{eq:multiscale-cost}.
For a pair in \eqref{eq:paired-frequencies}, nonnegativity of $\Lambda$
implies that both corresponding signals are at least $e^{-2.74E}$, and
\[
 \abs{\Delta[A(t)]\Delta[A(t')]}
 =U^2e^{-\Lambda(t)-\Lambda(t')}\ge e^{-2.74E}.
\]
Applying
\eqref{eq:evaluation-lower-bound} to each therefore gives
\begin{equation}\label{eq:pair-normalization}
 \abs{B_x(t)\overline{B_x(t')}}\ge e^{-2.75E}.
\end{equation}
For example, each evaluation is at least one third of its corresponding
signal once $E$ is large; the resulting factor $1/9$ is absorbed in $e^{-0.01E}$.

Take a monotone chain
$0=y_0,y_1,\ldots,y_J=\lambda$ with consecutive distances at most $b$ and
\[
 1\le J\le C(1+S/b).
\]
If $\lambda=0$, use just two endpoints and $J=1$.
For every interior $y_i$, choose a pair from
\eqref{eq:paired-frequencies} and define polynomials
\[
 V_0=1,\qquad
 V_i=\frac{A(t)\overline{A(t')}}{B_x(t)\overline{B_x(t')}}\quad(0<i<J),
 \qquad V_J=\frac{A(\lambda)}{B_x(\lambda)}.
\]
The denominators are nonzero by
\eqref{eq:endpoint-normalization}--\eqref{eq:pair-normalization}, and every
$V_i(x)$ equals one.  Regard these polynomials as vectors $v_i$ in the
Hilbert space induced by $\EE$.  They have degree at most $2k\le D$.

For $1\le i\le J$, the norm square of $v_i$ and its inner product with
$v_{i-1}$ are expectations of normalized balanced products: their signed
frequencies are respectively zero and
$y_i-y_{i-1}$, and each product has at most four factors.  The magnitude of its
normalizing denominator is at least $e^{-5.50E}$.  Our contradictory
assumption therefore gives
\[
 \abs{\norm{v_i}^2-1}\le\xi,\qquad
 \abs{\langle v_i,v_{i-1}\rangle-1}\le\xi,
 \qquad \xi=\epsilon e^{5.50E}=e^{-2.44E}.
\]
Since $v_0=1$ has norm one, Lemma~\ref{lem:gram-chain} yields
\[
 \abs{\Delta[A(\lambda)]}
 =\abs{B_x(\lambda)}\,
       \abs{\langle v_J,v_0\rangle-1}
 \le2UJ^2e^{-2.44E}<e^{-2.40E}.
\]
The final strict inequality follows by increasing $C_0$, because
$\log(2UJ^2)$ is less than $0.04E$ by \eqref{eq:multiscale-cost}.
This contradicts \eqref{eq:single-frequency} and proves the lemma.
\end{proof}

We have obtained a product of at most four windowed feature sums with signal
at least $e^{-7.94E}$ and total frequency at most $b$.  It remains to turn it
into one feature while keeping both its decay and frequency of order $S'$.
The Gaussian contour construction below supplies that final step of
Proposition~\ref{prop:multiscale}.
